\documentclass[11pt]{article}

\usepackage[a4paper,margin=2.4cm]{geometry}
\usepackage{fontspec}
\setmainfont{Helvetica Neue}
\usepackage{titlesec}
\usepackage{enumitem}
\usepackage{xcolor}
\usepackage{listings}
\usepackage{tabularx}
\usepackage{booktabs}
\usepackage{hyperref}
\usepackage{fancyhdr}
\usepackage{tcolorbox}
\tcbuselibrary{skins,breakable}

\definecolor{codebg}{RGB}{244,244,247}
\definecolor{codeframe}{RGB}{210,210,216}
\definecolor{accent}{RGB}{31,74,121}
\definecolor{warnbg}{RGB}{255,245,230}
\definecolor{warnframe}{RGB}{224,168,90}

\lstset{
  basicstyle=\ttfamily\small,
  breaklines=true,
  backgroundcolor=\color{codebg},
  frame=single,
  rulecolor=\color{codeframe},
  framesep=6pt,
  xleftmargin=2pt,
  xrightmargin=2pt,
  columns=fullflexible,
  keepspaces=true,
}

\newtcolorbox{note}{colback=codebg,colframe=codeframe,boxrule=0.5pt,arc=2pt,left=8pt,right=8pt,top=6pt,bottom=6pt,breakable}
\newtcolorbox{warn}{colback=warnbg,colframe=warnframe,boxrule=0.5pt,arc=2pt,left=8pt,right=8pt,top=6pt,bottom=6pt,breakable,title={Watch out},fonttitle=\bfseries\color{warnframe!60!black},coltitle=black}

\titleformat{\section}{\Large\bfseries\color{accent}}{\thesection}{0.6em}{}
\titleformat{\subsection}{\large\bfseries\color{accent}}{\thesubsection}{0.6em}{}

\pagestyle{fancy}
\fancyhf{}
\fancyhead[L]{\small\color{gray} Axion Strings AMR --- Cluster Guide}
\fancyhead[R]{\small\color{gray} \thepage}
\renewcommand{\headrulewidth}{0.3pt}

\hypersetup{
  colorlinks=true,
  linkcolor=accent,
  urlcolor=accent,
  citecolor=accent,
  pdftitle={Axion Strings AMR: Running on a Cluster},
  pdfauthor={Axion Strings AMR project}
}

\setlist[itemize]{leftmargin=1.4em}
\setlist[enumerate]{leftmargin=1.6em}

\title{\vspace{-1.5cm}\textbf{\Huge Axion Strings AMR}\\[0.3cm]\Large Getting Started on a Cluster}
\author{}
\date{\today}

\begin{document}
\maketitle
\vspace{-1.2cm}

\begin{note}
\textbf{Who this is for.} You can already \texttt{ssh} into a cluster and \texttt{scp}/\texttt{rsync} files
back and forth, but you don't need to be a systems person to follow this. Every command you need is given
in full --- copy, paste, adjust the few things flagged as ``change this''.
\end{note}

\tableofcontents

\section{What you're running}

This is a 3D simulation code for axion cosmic strings, built on the AMReX/GRTeclyn adaptive-mesh-refinement
(AMR) framework. ``Adaptive mesh'' means the simulation automatically uses a fine grid only where the
physics needs it (near strings) and a coarse grid everywhere else, which is what makes bigger, more
realistic runs affordable at all.

Practically, this means two things:
\begin{itemize}
  \item One compiled program (an \texttt{.ex} file) runs every kind of simulation this project does ---
        different physics comes from a plain-text \emph{parameter file}, not different code.
  \item A run writes two very different kinds of output: a handful of small text files with the physics
        summaries you actually want (string density, energy, spectra, \dots), and --- only if you ask for
        them --- large binary ``plotfiles'' with the full 3D field on disk. You will almost never want to
        move the plotfiles off the cluster; see Section~\ref{sec:getting-data-back}.
\end{itemize}

\section{Step 1: Get the code onto the cluster}

The repository uses two git \emph{submodules} (GRTeclyn and its dependency, AMReX) that must be checked
out alongside it. If the cluster's login node has internet access, the cleanest approach is to clone
directly there:

\begin{lstlisting}
ssh you@cluster.example.edu
git clone --recurse-submodules https://github.com/<your-username>/axion_strings_adaptive.git
cd axion_strings_adaptive
\end{lstlisting}

If the cluster has no internet access from compute/login nodes (common on HPC systems), clone and prepare
the submodules on your own machine first, then copy the whole tree over with \texttt{rsync} (much faster
than \texttt{scp} for a big directory tree, and it can resume if interrupted):

\begin{lstlisting}
# on your own machine
git clone --recurse-submodules https://github.com/<your-username>/axion_strings_adaptive.git
rsync -avz --progress axion_strings_adaptive/ you@cluster.example.edu:~/axion_strings_adaptive/
\end{lstlisting}

\begin{note}
Do \emph{not} copy a compiled executable from your laptop to the cluster --- it was built for a different
processor and operating system and will not run. Always build fresh on the cluster (Step 2).
\end{note}

\section{Step 2: Build it on the cluster}

\subsection{Load a compiler and MPI}

Clusters make software available through \texttt{module} commands rather than having everything installed
system-wide. The exact module names differ between clusters, so start by looking:

\begin{lstlisting}
module avail mpi
module avail fftw
\end{lstlisting}

On DESY's Maxwell cluster (the system this guide's SLURM template was written against) this looks like:

\begin{lstlisting}
module load mpi/openmpi-x86_64
module load fftw
\end{lstlisting}

\subsection{Tell the build where FFTW is}

The code uses FFTW (for generating initial conditions), and the build needs to know where it's installed.
After loading the \texttt{fftw} module, find its location:

\begin{lstlisting}
module show fftw
\end{lstlisting}

Look for a line setting something like \texttt{FFTW\_ROOT}, \texttt{EBROOTFFTW}, or a \texttt{-I.../include}
flag --- that directory (the one \emph{containing} \texttt{include/} and \texttt{lib/}) is what you pass as
\texttt{FFTW\_HOME} below.

\subsection{Build}

This guide builds with multi-threading (OpenMP) switched on --- each MPI rank uses several CPUs at once,
not just one --- since that's the setup you actually want for a real allocation:

\begin{lstlisting}
cd axion_strings_adaptive/AxionStrings
make COMP=gnu FFTW_HOME=/path/from/module/show USE_OMP=TRUE -j8
\end{lstlisting}

\begin{itemize}
  \item \texttt{COMP=gnu} tells it to use GCC, the standard choice on Linux clusters (on the project's own
        Mac, \texttt{COMP=llvm} is used instead --- don't copy that part of any Mac instructions you see).
  \item \texttt{USE\_OMP=TRUE} is what turns on multi-threading --- see Section~\ref{sec:openmp} for what
        this means for how you fill in the SLURM script. The bulk of the per-step cost (the physics update
        itself, and every energy/string-count diagnostic) is already written in a form AMReX itself splits
        across threads automatically when built this way --- no further code changes needed to benefit from
        it.
  \item \texttt{-j8} builds with 8 parallel compile jobs; raise it if you have more cores available on the
        login node, but check the cluster's policy on running CPU-heavy work on login nodes first --- if
        it's restricted, submit the build itself as a short batch job (Section~\ref{sec:submitting}, but
        with the \texttt{make} command in place of \texttt{mpirun}, and no MPI needed for the build). This
        \texttt{-j8} is about compiling the code faster and is unrelated to \texttt{USE\_OMP} (which is
        about the compiled program's own later multi-threading).
\end{itemize}

A successful build takes a few minutes and ends with an executable named something like
\texttt{AxionStrings3d.gnu.FLOAT.MPI.OMP.ex} sitting in \texttt{AxionStrings/} --- note the \texttt{.OMP}
tag, confirming multi-threading is compiled in. If it fails, the very last handful of lines of the error is
almost always a missing module (wrong \texttt{FFTW\_HOME}, or an MPI compiler wrapper not on \texttt{PATH})
--- re-check Steps 2.1--2.2 before anything else.

\subsection{Ranks and threads: how to split up your CPUs}
\label{sec:openmp}

You now have two independent dials for spreading a run across CPUs: \texttt{--ntasks} (MPI ranks, each a
separate process with its own chunk of the domain) and \texttt{--cpus-per-task} (OpenMP threads, each rank
splitting its own work further across several CPUs at once). Their product is the total number of CPUs the
job reserves. A reasonable starting point is a handful of threads per rank (\texttt{--cpus-per-task=4} is
used in the template below) rather than all-ranks-one-thread-each or one-rank-all-threads --- adjust to
taste once you have a feel for how your specific runs scale.

\begin{warn}
A \texttt{USE\_OMP=TRUE} build does nothing by itself until you also tell it how many threads to use, via
the standard \texttt{OMP\_NUM\_THREADS} environment variable --- \textbf{always set it to match
\texttt{\$SLURM\_CPUS\_PER\_TASK}} in your submission script (the template in Section~\ref{sec:submitting}
already does this --- keep that line if you copy it). Leaving it unset is not merely a missed optimisation:
OpenMP then defaults to using \emph{every} core it can see on the node, which badly oversubscribes the
machine as soon as more than one MPI rank shares that node (the normal case). \texttt{--cpus-per-task} and
\texttt{OMP\_NUM\_THREADS} always travel together --- never set one without the other.
\end{warn}

\section{Step 3: Organise your files on the cluster}
\label{sec:organise}

Most clusters give you (at least) two separate storage areas with very different rules:

\begin{tabularx}{\textwidth}{@{}>{\raggedright\arraybackslash}p{3.1cm} X@{}}
\toprule
\textbf{Area} & \textbf{Use it for} \\
\midrule
\texttt{\$HOME} & The code itself (this repo). Usually small quota, backed up, \emph{not} meant for run
output. \\[0.15cm]
\texttt{\$SCRATCH} / \texttt{\$WORK} (name varies) & Everything a simulation writes. Usually large quota,
\emph{not} backed up, and some clusters auto-delete old files here --- treat it as temporary. \\
\bottomrule
\end{tabularx}

\vspace{0.3cm}
Run \texttt{ls -la \$HOME} and check your cluster's documentation (or ask IT) for the scratch area's
actual environment variable if \texttt{\$SCRATCH} isn't it.

\textbf{Recommended layout} --- one directory per run, never reused, so nothing from an old run silently
mixes with a new one:

\begin{lstlisting}
$SCRATCH/axion_runs/
    2026-09-21_N512_seed1/
        params.txt          # the exact parameter file used
        slurm.sh             # the exact submission script used
        run.log               # this run's console output
        network_scalars.dat   # small diagnostics (keep)
        axion_spectrum.dat     # small diagnostics (keep)
        axion_projection.dat    # small diagnostics (keep, if enabled)
        plt00000/  plt00100/  ...  # plotfiles, if amr.plot_int is set (usually large -- see Sec. 9)
        chk00000/                    # checkpoints, if amr.check_int is set
\end{lstlisting}

Copying \texttt{params.txt} and \texttt{slurm.sh} into the run's own directory (rather than only keeping
one shared copy) means that six months from now you can look at any run directory and know exactly what
produced it, with no guessing from memory or from a shared file that's since been edited.

\section{Step 4: Set up a run --- the parameter file}

Every run is controlled by one plain-text file (e.g.\ \texttt{params.txt}), lines of the form
\texttt{key = value}. Existing examples live in \texttt{AxionStrings/params\_*.txt} --- copying the closest
existing one and editing it is usually easier than writing one from scratch.

\subsection{The parameters you'll actually touch}

\begin{description}[style=nextline,leftmargin=1.2em,itemsep=0.4em]
\item[\texttt{axion\_strings.c0}]
  Which physical system: \texttt{0} = physical strings (the real thing), \texttt{1} = ``fat'' string
  (cheaper, approximate), \texttt{2} = ``Moore'' mode. Start with \texttt{0} unless told otherwise.
\item[\texttt{axion\_strings.N}]
  Overall box size / resolution target. This is the single biggest lever on both physical realism and
  cost --- see Section~\ref{sec:scaling-up}.
\item[\texttt{axion\_strings.N1}, \texttt{N2}]
  Leave these at \texttt{1.0} unless someone has specifically told you otherwise --- they fix two resolution
  targets (``horizon fits in the box'' and ``string core is resolved'') that the project has already
  settled on.
\item[\texttt{axion\_strings.log\_mr\_over\_h\_i}]
  How far into the simulation's natural time range the run \emph{starts}. Bigger box (\texttt{N}) sizes
  can afford to start later (a bigger number) without the string network being unrealistically fresh.
  \texttt{2.0} is a safe, already-used starting point.
\item[\texttt{amr.max\_level}]
  How many levels of automatic refinement to allow (\texttt{0} = no refinement, a plain uniform grid).
  Higher is more realistic \emph{and} usually much cheaper for the same finest resolution (adaptive
  refinement only pays for fine cells where they're needed) --- see the timing note in
  Section~\ref{sec:scaling-up}. \texttt{2} is a good, validated starting point.
\item[\texttt{axion\_strings.ic\_mode}]
  How the simulation starts. \texttt{fourier\_relaxed} (a randomised starting field, settled into a
  realistic tangle before the ``real'' run begins) is what production runs use --- it needs the four
  companion parameters just below, all of which are \emph{required} (the run refuses to start without
  them, there is no built-in fallback).
\item[\texttt{axion\_strings.pre\_evolution.seed}]
  The random seed for that starting field. \textbf{This is your ensemble knob} --- run the same parameter
  file with different seeds to get independent realisations of the same physics (see the SLURM array
  example in Section~\ref{sec:submitting}).
\item[\texttt{amr.plot\_int}]
  How often (in steps) to write a full 3D plotfile. These are large. Leave at \texttt{-1} (off) unless you
  specifically want to visualise the field; the physics diagnostics you actually analyse
  (\texttt{network\_scalars.dat} etc.) are controlled separately below and are cheap.
\item[\texttt{axion\_strings.output\_first\_log\_mr\_over\_h}]
  Where in the run to start writing rows to the small diagnostics files. There's no built-in default ---
  in practice this is almost always set equal to \texttt{log\_mr\_over\_h\_i} just above (start recording
  from the very beginning of the main run).
\item[\texttt{axion\_strings.output\_delta\_log\_mr\_over\_h}]
  How often, after that, to write another row. \texttt{0.1--0.2} is a typical spacing.
\end{description}

\subsection{The \texttt{fourier\_relaxed} companion parameters}
\label{sec:fourier-relaxed-params}

These four are required whenever \texttt{axion\_strings.ic\_mode = fourier\_relaxed} (the recommended,
default choice above). All four have settled, essentially fixed typical values, or --- for \texttt{gamma}
--- are simply left unset so the program derives it for you.

\begin{description}[style=nextline,leftmargin=1.2em,itemsep=0.4em]
\item[\texttt{axion\_strings.pre\_evolution.gamma}]
  \textbf{Leave this unset.} The program derives it automatically as \texttt{1/dx\_base} (the coarsest-level
  grid spacing) and prints what it derived as part of its startup box plan --- this is the resolution the
  relaxation phase needs to run at so the string cores it produces are resolved consistently with what the
  main run's own grid looks like the moment it starts (relies on \texttt{log\_mr\_over\_h\_i} being below
  the level-1 threshold printed just above it, i.e.\ the main run genuinely starting on the coarsest grid
  alone --- the program checks this itself and refuses to start if it isn't true, rather than silently
  relaxing at the wrong resolution). You can still set it by hand if you have a specific reason to (e.g.\
  experimenting with a different relaxation resolution) --- the program then checks your value against what
  it would have derived and prints a warning (not an error --- the run still proceeds with your value) if
  they disagree, so an accidental stale value from a previous box size doesn't go unnoticed.
\item[\texttt{axion\_strings.pre\_evolution.xi\_target}]
  Stop relaxing once the (measured) string density first drops to this value. \textbf{Typical value:
  \texttt{1.0}} (used throughout the project's own validated examples; some production runs instead target
  a lower value like \texttt{0.7} for a more settled starting tangle at the cost of a longer relaxation
  phase).
\item[\texttt{axion\_strings.pre\_evolution.k\_max\_over\_mr}]
  The highest wavenumber (in units of $m_r$) populated in the random starting field. \textbf{Typical value:
  \texttt{64}}, used in every validated example at \texttt{amr.max\_level=2}. If you're testing a
  \emph{much} coarser base grid or a deeper hierarchy than those examples, this is worth being cautious
  with --- \texttt{docs/STATUS.md} documents a specific pathology where too high a value here (relative to
  how coarse the base grid is) over-tags the whole domain for refinement; if a test run's very first regrid
  tags close to 100\% of the domain unexpectedly, this is the first parameter to suspect.
\item[\texttt{axion\_strings.pre\_evolution.mean\_square\_variance}]
  How much random power to start with. \textbf{Always \texttt{0.5}} in this project's convention (chosen
  so the field's total mean-square amplitude works out to the natural unit value) --- not something to
  change casually.
\end{description}

\subsection{A note on \texttt{amr.max\_grid\_size}}

This one doesn't change the physics at all, only how the domain is chopped up for parallel processing ---
but it needs to roughly match how many MPI ranks you plan to run with (Section~\ref{sec:submitting}). A
smaller \texttt{amr.max\_grid\_size} makes more, smaller chunks, letting you usefully spread the coarse
level across more ranks. As a rule of thumb: (grid size at that level) $\div$ \texttt{max\_grid\_size},
cubed, should be at least as big as your \emph{rank} count --- specifically ranks, not total CPUs: each
rank's own threads (Section~\ref{sec:openmp}) already split the cells \emph{within} whatever chunks that
rank owns, so this box count doesn't need to grow with \texttt{--cpus-per-task} the way it would need to
grow with \texttt{--ntasks}.

\subsection{Scaling up to a bigger test}
\label{sec:scaling-up}

\begin{enumerate}
  \item Take the closest existing \texttt{params\_*.txt} as a starting template (e.g.\
        \texttt{params\_full\_test\_640.txt}).
  \item Raise \texttt{axion\_strings.N}. Leave \texttt{N1}/\texttt{N2} alone.
  \item If your template has an \texttt{axion\_strings.pre\_evolution.gamma} line (every existing example
        does), delete it --- it was computed for the \emph{old} box size and won't match the new one, and
        while the program will still run and only warn about the mismatch (Section~\ref{sec:fourier-relaxed-params}),
        the relaxation phase would silently run at the wrong resolution if you missed that warning. Leaving
        \texttt{gamma} out entirely sidesteps the issue and lets the program derive the right value for you.
  \item \textbf{Always do a very short test run first} (see the box below) before committing a multi-hour
        job to the queue.
\end{enumerate}

\begin{warn}
The program prints its own ``box plan'' (grid spacing, box size, refinement schedule) in the first few
lines of output, every time it starts, \emph{before} it starts stepping. Always read those first few lines
after raising \texttt{N} and check nothing looks surprising (in particular, an unexpectedly coarse grid
spacing has previously caused a run to tag far more of the domain for refinement than intended, quietly
making a run much more expensive than planned). A 2--3 minute local test with a handful of steps
(\texttt{evolution.max\_steps=50} added on the command line) is enough to see this printout and catch a
misconfiguration before it costs a multi-hour cluster allocation.
\end{warn}

\subsection{A worked example: a moderately big cluster run}
\label{sec:worked-example}

\texttt{AxionStrings/params\_cluster\_512base\_2level.txt} is a complete, ready-to-run example at a scale
genuinely worth a cluster (not just a laptop test): a $512^3$ base grid with 2 refinement levels
($N_\mathrm{effective}=2048$), physical strings, every optional diagnostic in Section~\ref{sec:advanced}'s
``Extra diagnostics'' turned on so you can see all of them wired up in one place, and the regrid-buffer
mechanism from Section~\ref{sec:regrid-buffer} used instead of hand-set values. Its own header comment shows
the full box-plan arithmetic worked by hand, the same way this section asks you to check it --- and it has
actually been smoke-tested (a few steps, confirmed no crash and an exact match between the hand-worked
numbers and what the program itself printed) before being included here, not just written and assumed
correct.

Use it directly, or --- more usefully --- as the template Step 1 above refers to: copy it, change
\texttt{axion\_strings.N} for your own target size, delete the (still-present, now stale) comment block
doing the old box-plan arithmetic by hand, and let the short test run's own printout give you the new
numbers instead.

\section{Less commonly changed options}
\label{sec:advanced}

Everything in this section has a sensible default and most runs never need to touch it --- skip it on a
first read and come back if you have a specific reason to (a run behaving oddly, or someone asking you to
try a specific value).

\subsection{How often refinement is re-checked, and its safety buffer}
\label{sec:regrid-buffer}

The adaptive mesh doesn't re-decide where to refine on every single step (that would be wasteful) --- it
re-checks periodically, then grows a safety buffer of extra-refined cells around every flagged region so
that nothing genuinely fast-moving can slip out of the fine grid in between checks.

\textbf{There are exactly two parameters to set here, together, and nothing else in this subsection needs
touching once you have:}

\begin{description}[style=nextline,leftmargin=1.2em,itemsep=0.4em]
\item[\texttt{axion\_strings.tagging.regrid\_interval\_steps}]
  How often, in steps, the simulation re-decides where to refine. \textbf{Typical value: \texttt{2}}
  (matching every one of the project's own validated example runs). Smaller tracks the strings more
  tightly (fewer stale/wasted fine cells) but costs more time spent re-deciding; larger is cheaper
  bookkeeping but requires a wider safety buffer (derived automatically below), meaning more volume held at
  finer resolution than strictly necessary.
\item[\texttt{axion\_strings.tagging.buffer\_safety\_factor}]
  A multiplier $\geq 1.0$ on top of the bare-minimum buffer width the code derives from
  \texttt{regrid\_interval\_steps} above. \textbf{Default: \texttt{1.0}} (the bare minimum) --- raise it
  (e.g.\ to \texttt{1.5} or \texttt{2.0}) only for extra headroom when trying a new, not-yet-trusted
  combination of parameters, at the cost of refining a bit more volume than the bare minimum.
\end{description}

\begin{warn}
Setting these two is an \emph{alternative} to setting \texttt{amr.regrid\_int}/\texttt{amr.n\_error\_buf}
yourself, not an addition to it --- once \texttt{regrid\_interval\_steps} is set, the run derives both of
those for you automatically and prints what it derived (and why) in its startup output, right alongside the
box plan (worth reading alongside the box-plan check in Section~\ref{sec:scaling-up}'s warning box).
\textbf{Don't also set \texttt{amr.regrid\_int}/\texttt{amr.n\_error\_buf} by hand in the same parameter
file} --- if you do, the run cross-checks them against what it derived and stops with an error rather than
silently picking one. The rest of this box explains what those two lower-level parameters mean, only for
the (uncommon) case where you want to bypass the mechanism above entirely and set them directly instead:

\begin{description}[style=nextline,leftmargin=1.2em,itemsep=0.3em]
\item[\texttt{amr.regrid\_int}] Steps between refinement re-checks, set directly. \texttt{-1} on a single
  level disables regridding entirely (used by the wave-packet boundary tests in
  \texttt{AxionStrings/analysis/wave\_boundary\_test/} to hold a grid fixed).
\item[\texttt{amr.n\_error\_buf}] The safety buffer width itself, in cells, set directly.
\end{description}
\end{warn}

\subsection{How sensitive the refinement trigger is (off by default)}

What actually finds strings for refinement, for context on the two optional parameters below: the tagger's
\emph{primary} criterion is an exact, parameter-free topological test --- does the field's phase wind by a
full $\pm 2\pi$ loop around a cell's corner plaquette? That's an integer (a string pierces this cell or it
doesn't), not a threshold on a continuous quantity, so it has no ``units'' question and nothing for you to
set --- it's always on. \emph{Which levels are even allowed to exist yet} is controlled separately, by the
schedule (\texttt{log\_mr\_over\_h\_i} and the box-plan-derived level-addition thresholds); the two
parameters below don't touch that either. Both default to effectively infinite (i.e.\ off) and were found,
in earlier testing, to flag essentially the entire domain if turned on carelessly --- leave them alone
unless you have a specific reason to experiment with gradient-based tagging, and if you do, read
\texttt{StringTagger.hpp}'s own header comment first for the full reasoning.

\begin{description}[style=nextline,leftmargin=1.2em,itemsep=0.4em]
\item[\texttt{axion\_strings.tagging.gradient\_threshold}]
  A \emph{second}, independent, optional trigger --- despite the name, actually a \emph{Laplacian}
  (curvature) test, applied separately to each of the field's two real components ($\psi_1$, $\psi_2$; not
  a decomposed ``axion'' or ``phase'' quantity specifically --- that split is what the next parameter does).
  A cell is tagged if $dx^2\,|\nabla^2\psi_i|$ exceeds this threshold, for either component.
  \textbf{Units: dimensionless} --- $dx^2$ (comoving length$^2$) exactly cancels the Laplacian's
  $1/\mathrm{length}^2$, leaving a pure number of order ``how much $\psi$ curves across one cell, relative to
  its own natural $O(1)$ scale'' ($\psi$ is normalised by the vacuum value $v=1$ in this code's units).
  Buschmann et al.\ 2108.05368's own value is \texttt{0.04}, but this project's own testing found that value
  tags essentially the whole domain for its field statistics --- not a starting point to trust as-is.
\item[\texttt{axion\_strings.tagging.radial\_gradient\_threshold}]
  A \emph{third}, independent, optional trigger, this one genuinely radial: $r=|\psi|$, the field's
  amplitude alone (properly decomposed away from the phase direction, not just combining
  $\nabla\psi_1$/$\nabla\psi_2$ naively). A cell is tagged if $dx\,|\nabla r|$ exceeds this threshold.
  \textbf{Units: also dimensionless} (same cancellation, one power of $dx$ against $\nabla r$'s
  $1/\mathrm{length}$) --- but this is a \emph{first} derivative of a different, derived quantity, not the
  same operation as \texttt{gradient\_threshold}'s Laplacian of the raw components. Both end up
  dimensionless on the same rough scale (``per-cell change in units of the field's own $O(1)$ scale''), but
  a given numeric value doesn't mean the same thing for both --- don't assume tuning one tells you anything
  about the other.
\end{description}

\subsection{Domain decomposition, the other half}

\begin{description}[style=nextline,leftmargin=1.2em,itemsep=0.4em]
\item[\texttt{amr.blocking\_factor}]
  Every box's size (in each dimension) must be a multiple of this. Works alongside
  \texttt{amr.max\_grid\_size} (Section~\ref{sec:scaling-up}) to control how the domain is chopped up for
  parallel ranks --- the project's own example files use \texttt{16} or \texttt{32}; there's rarely a
  reason to deviate unless you're chasing a specific load-balancing problem.
\end{description}

\subsection{Timestep safety margin}

\begin{description}[style=nextline,leftmargin=1.2em,itemsep=0.4em]
\item[\texttt{evolution.dt\_multiplier}]
  The CFL safety factor: how large a timestep to take as a fraction of the finest active grid spacing.
  \texttt{0.2} is the project's standard, conservative choice. Raising it makes runs faster but risks
  numerical instability --- not something to change without a specific reason and a stability check.
\end{description}

\subsection{Cosmological expansion rate}

\begin{description}[style=nextline,leftmargin=1.2em,itemsep=0.4em]
\item[\texttt{axion\_strings.a\_inv}]
  Which cosmological era the box expands through (\texttt{2.0} = radiation domination, the project's
  standard choice throughout). Only change this if you specifically intend to simulate a different era ---
  it changes the box-planning formulas, not just a cosmetic label.
\end{description}

\subsection{Extra diagnostics}

These add real cost (the spectrum needs an FFT of the whole level-0 grid; projections write an extra file
per snapshot), which is why they default off.

\begin{description}[style=nextline,leftmargin=1.2em,itemsep=0.4em]
\item[\texttt{axion\_strings.masking.scheme}, \texttt{axion\_strings.masking.threshold}]
  \texttt{scheme} defaults to \texttt{A} (unscreened) if left unset. Set it to \texttt{B} (with a
  \texttt{threshold} in $(0,1)$ --- \textbf{required} in that case, no default) to additionally get
  \emph{screened} energy/tension numbers --- ``screened'' excludes a region around each string core from
  the energy sum, which some analyses want and others don't. Values \texttt{0.8}, \texttt{0.9} and
  \texttt{0.95} have all been used by this project at different times without settling on one, so treat any
  particular value as a choice to state explicitly in your own notes, not an established default.
  \texttt{axion\_strings.compute\_spectrum} (below) requires scheme B.
\item[\texttt{axion\_strings.compute\_spectrum}]
  Set to \texttt{1} to write \texttt{axion\_spectrum.dat} (needed for the spectrum panel of
  Section~\ref{sec:analyse-locally}'s standard plots). Requires \texttt{axion\_strings.masking.scheme = B}
  (just above).
\item[\texttt{axion\_strings.save\_projection}]
  Set to \texttt{1} to additionally save a 2D projection of the field at each diagnostic snapshot, useful
  for visualising the string network's evolution over time rather than only its summary statistics.
\end{description}

\subsection{Checkpointing and restart}

\begin{description}[style=nextline,leftmargin=1.2em,itemsep=0.4em]
\item[\texttt{amr.check\_int}]
  How often (in steps) to write a full checkpoint you can restart from. \texttt{-1} (the default in every
  example file) turns it off. Turn it on for any run that might not finish inside one job's walltime --- see
  the warning box in Section~\ref{sec:watch-it-run}. Restarting from a checkpoint changes how you invoke the
  executable; ask if you need this and it isn't already covered by an example you've found.
\end{description}

\section{Step 5: Write and submit the SLURM script}
\label{sec:submitting}

Below is a template adapted from an existing working example for DESY's Maxwell cluster. If you're on a
different cluster, the \texttt{--partition}, module names, and \texttt{OMPI\_MCA\_btl} line are the parts
most likely to need changing --- ask your cluster's documentation or support for the local equivalents.

\begin{lstlisting}
#!/bin/bash -l
#SBATCH --job-name=axion_test
#SBATCH --output=%x_%A_%a.log
#SBATCH --array=1-4                 # one job per seed: 4 independent realisations
#SBATCH -N 1                        # nodes
#SBATCH --ntasks=8                  # MPI ranks (see amr.max_grid_size note, Sec. 5.3)
#SBATCH --cpus-per-task=4           # OpenMP threads per rank, Sec. 3.4 -- 8x4 = 32 CPUs total
#SBATCH --time=04:00:00             # walltime: raise for bigger/longer runs
#SBATCH --partition=maxcpu

module load mpi/openmpi-x86_64
export OMPI_MCA_btl=^vader,tcp,openib
export OMP_NUM_THREADS=$SLURM_CPUS_PER_TASK   # must match --cpus-per-task, Sec. 3.4

CODE_DIR=$HOME/axion_strings_adaptive/AxionStrings
RUN_DIR=$SCRATCH/axion_runs/$(date +%F)_test_seed${SLURM_ARRAY_TASK_ID}
mkdir -p "$RUN_DIR"
cp "$CODE_DIR/AxionStrings3d.gnu.FLOAT.MPI.OMP.ex" "$CODE_DIR/params.txt" "$RUN_DIR/"
cd "$RUN_DIR"
cp "$0" slurm.sh   # keep a copy of the exact script that produced this run

mpirun -np 8 ./AxionStrings3d.gnu.FLOAT.MPI.OMP.ex params.txt \
    axion_strings.pre_evolution.seed=${SLURM_ARRAY_TASK_ID} \
    --bind-to none \
    > run.log 2>&1
\end{lstlisting}

A few things worth calling out explicitly:
\begin{itemize}
  \item \textbf{\texttt{\#SBATCH --array=1-4}} is what runs 4 independent realisations (different random
        seeds) as one job submission --- each array element gets its own
        \texttt{\$SLURM\_ARRAY\_TASK\_ID}, which the script both uses as the seed \emph{and} as part of the
        run directory name, so the four never collide.
  \item Parameters after the parameter-file name on the \texttt{mpirun} line
        (\texttt{axion\_strings.pre\_evolution.seed=\dots}) \emph{override} whatever's in \texttt{params.txt}
        --- handy for exactly this kind of per-job tweak without needing four separate parameter files.
  \item \textbf{\texttt{--bind-to none}} matters more here than it would without threads: it stops Open MPI
        pinning each rank to a single CPU, which would otherwise trap that rank's 4 threads fighting over
        one core instead of spreading across the 4 \texttt{--cpus-per-task} reserved for it.
  \item \texttt{--ntasks} is the number of MPI ranks --- it should match \texttt{-np} on the
        \texttt{mpirun} line. \texttt{--cpus-per-task} is threads \emph{per rank}, and only does anything
        useful together with a \texttt{USE\_OMP=TRUE} build and a matching \texttt{OMP\_NUM\_THREADS}
        (Section~\ref{sec:openmp}, both already in this template) --- if you ever build \emph{without}
        \texttt{USE\_OMP=TRUE}, set \texttt{--cpus-per-task} back to \texttt{1}, or the extra CPUs just sit
        idle.
  \item Everything writes into a fresh, timestamped run directory under \texttt{\$SCRATCH}, matching
        Section~\ref{sec:organise}.
\end{itemize}

Submit it, and check it's queued:

\begin{lstlisting}
sbatch slurm.sh
squeue -u $USER
\end{lstlisting}

\section{Step 6: Watch it run}
\label{sec:watch-it-run}

Once a job starts, its log file (\texttt{run.log}, from the redirection in the script above) fills in
real time:

\begin{lstlisting}
tail -f $SCRATCH/axion_runs/<run directory>/run.log
\end{lstlisting}

A healthy run shows repeating blocks like:

\begin{lstlisting}
[Level 0 step 214] ADVANCE at time 3.02 with dt = 0.0142
[Level 0 step 214] average evolution speed = 812.4 code units/h
[Level 0 step 214] Advanced 2097152 cells
\end{lstlisting}

\texttt{average evolution speed} is the code's own built-in throughput readout --- if it's dropping sharply
and staying low, the simulation has likely started refining a large fraction of the domain (check the
\texttt{REGRID} lines just above for how much). \texttt{network\_scalars.dat} (in the same run directory)
fills in alongside it with the physics summaries and can be watched the same way (\texttt{tail -f}) for a
terser, physics-only view.

\begin{warn}
If a job hits its \texttt{--time} limit before finishing, SLURM kills it --- it does not pause and resume
automatically. If you expect a run might need longer than one job's walltime, turn on checkpointing
(\texttt{amr.check\_int}, an interval in steps) so you can resubmit and continue from the last checkpoint
rather than losing the run. This is not covered further here --- ask if you need it, since it changes the
submission script slightly (the executable needs to be told which checkpoint to restart from).
\end{warn}

\section{Step 7: Bring the results home}
\label{sec:getting-data-back}

You almost always want only the small diagnostic files, not the plotfiles (each plotfile is the entire 3D
field on disk and can be many gigabytes; the diagnostics already capture everything the standard analysis
plots need). From your own machine:

\begin{lstlisting}
rsync -avz \
    you@cluster.example.edu:$SCRATCH/axion_runs/<run directory>/network_scalars.dat \
    you@cluster.example.edu:$SCRATCH/axion_runs/<run directory>/axion_spectrum.dat \
    you@cluster.example.edu:$SCRATCH/axion_runs/<run directory>/axion_projection.dat \
    ~/axion_runs/<run directory>/
\end{lstlisting}

(Drop any of the three filenames that weren't produced by your particular run --- e.g.\
\texttt{axion\_spectrum.dat} only appears if \texttt{axion\_strings.compute\_spectrum=1} was set.) It's
worth also grabbing \texttt{run.log}, \texttt{params.txt} and \texttt{slurm.sh} from the same directory ---
tiny, and they're the record of exactly what produced the data.

\section{Step 8: Analyse it locally}
\label{sec:analyse-locally}

The repository's \texttt{AxionStrings/analysis/} scripts take those diagnostic files directly and produce
the project's standard plots. On your own machine, with the repository's Python environment active
(see the analysis folder's own \texttt{README.md} for setup if you haven't used it before):

\begin{lstlisting}
cd axion_strings_adaptive/AxionStrings/analysis
python3 plot_run_summary.py \
    ~/axion_runs/<run directory>/network_scalars.dat \
    ~/axion_runs/<run directory>/axion_spectrum.dat \
    ~/axion_runs/<run directory>/plots_out \
    <a_inv> <c0> <N> [optional title tag]
\end{lstlisting}

This produces the standard set: string density $\xi$, energies (raw and normalised), string tension, and
the axion spectrum, both raw and normalised. Fill in \texttt{<a\_inv>}, \texttt{<c0>}, \texttt{<N>} with the
exact values from the \texttt{params.txt} you ran with (they affect the unit conversions, so they need to
match, not just be ``close'').

If you don't have an \texttt{axion\_spectrum.dat} (i.e.\ you didn't turn spectrum output on), pass
\texttt{None} in its place --- the script skips the spectrum panel and still produces the rest.

\section{Quick reference}

\begin{tabularx}{\textwidth}{@{}>{\raggedright\arraybackslash}p{3.6cm} X@{}}
\toprule
\textbf{Task} & \textbf{Command} \\
\midrule
Get code onto cluster & \texttt{git clone --recurse-submodules \dots} or \texttt{rsync -avz \dots} \\
Build & \texttt{make COMP=gnu FFTW\_HOME=\dots USE\_OMP=TRUE -j8} \\
Submit a run & \texttt{sbatch slurm.sh} \\
Check queue & \texttt{squeue -u \$USER} \\
Watch a running job & \texttt{tail -f run.log} \\
Cancel a job & \texttt{scancel <jobid>} \\
Bring results home & \texttt{rsync -avz user@cluster:\dots/*.dat ~/axion\_runs/\dots/} \\
Make the standard plots & \texttt{python3 plot\_run\_summary.py \dots} \\
\bottomrule
\end{tabularx}

\vspace{0.5cm}
\begin{note}
If something in this guide doesn't match what you see (a module name, a file that isn't there, an error
message), that's much more likely to be a cluster-specific detail this guide couldn't anticipate than a
mistake on your end --- check the cluster's own documentation for the specific piece (module names and
partition names are the most commonly different), and the rest of the workflow should carry over unchanged.
\end{note}

\end{document}
